
\documentclass{article}
\usepackage{colortbl}
\usepackage{makecell}
\usepackage{multirow}
\usepackage{supertabular}

\begin{document}

\newcounter{utterance}

\centering \large Interaction Transcript for game `dond', experiment `semi\_en', episode 8 with qwen3.5{-}9b.
\vspace{24pt}

{ \footnotesize  \setcounter{utterance}{1}
\setlength{\tabcolsep}{0pt}
\begin{supertabular}{c@{$\;$}|p{.15\linewidth}@{}p{.15\linewidth}p{.15\linewidth}p{.15\linewidth}p{.15\linewidth}p{.15\linewidth}}
    \# & $\;$A & \multicolumn{4}{c}{Game Master} & $\;\:$B\\
    \hline

    \theutterance \stepcounter{utterance}  
    & & \multicolumn{4}{p{0.6\linewidth}}{
        \cellcolor[rgb]{0.9,0.9,0.9}{
            \makecell[{{p{\linewidth}}}]{
                \texttt{\tiny{[P1$\langle$GM]}}
                \texttt{You are playing a negotiation game in which you have to agree on how to divide a set of items among you and another player.} \\
\\ 
\texttt{Rules:} \\
\texttt{(a) You and the other player are given a set of items. Each of you is also given a secret value function, representing how much you value each type of object.} \\
\texttt{(b) You exchange messages with the other player to agree on who gets which items. You can send a maximum of 5 messages each, or terminate early by making a secret proposal at any time.} \\
\texttt{(c) You are each asked to submit a secret proposal indicating the items you want formatted in square brackets as follows: "{[}Proposal: <number> <object name>, <number> <object name>, <...>{]}"} \\
\texttt{(d) If your proposals are complementary, i.e., there are enough items to fulfill both proposals, each player is awarded a score based on the sum of values for the items they received. Otherwise, both of you get zero points.} \\
\texttt{(e) Your goal is to maximize the score you receive.} \\
\\ 
\texttt{Let us start.} \\
\\ 
\texttt{The set of available items is:} \\
\\ 
\texttt{1 phone, 2 leaves, 2 oranges} \\
\\ 
\texttt{Your secret value function is:} \\
\\ 
\texttt{phone: 4, leaf: 0, orange: 3} \\
\\ 
\texttt{IMPORTANT: Your messages, unless it is the secret proposal, are directly transmitted to the other player, so do not include any response to the rules or text announcing your message. To make a secret proposal, use the indicated format. Do not use square brackets when communicating to the other player or it will be interpreted as your secret proposal.} \\
            }
        }
    }
    & \\ \\

    \theutterance \stepcounter{utterance}  
    & \multicolumn{4}{p{0.6\linewidth}}{
        \cellcolor[rgb]{0.9,0.9,0.9}{
            \makecell[{{p{\linewidth}}}]{
                \texttt{\tiny{[P1$\rangle$GM]}}
                \texttt{I value the phone at 4 and oranges at 3, while leaves are worthless to me (0). The other player's values are unknown. To maximize my score, I need to secure the phone and at least one orange. Since I don't know what the other player wants, I should propose a bundle that includes the phone and an orange, hoping they accept a deal where I take those items and they take the remaining items (the other orange and the leaves). If they want the phone, they will reject this, and we both get zero, but that's the risk of not knowing their preferences. A safer initial move is to propose a specific allocation that gives me the high{-}value items I need.} \\
\\ 
\texttt{Let's propose: I take the phone and one orange. This leaves the other orange and the leaves for them.} \\
\\ 
\texttt{Proposal: 1 phone, 1 orange} \\
            }
        }
    }
    & & \\ \\

    \theutterance \stepcounter{utterance}  
    & & & \multicolumn{2}{p{0.3\linewidth}}{
        \cellcolor[rgb]{0.9,0.9,0.9}{
            \makecell[{{p{\linewidth}}}]{
                \texttt{\tiny{[GM$|$GM]}}
                \texttt{continue} \\
            }
        }
    }
    & & \\ \\

    \theutterance \stepcounter{utterance}  
    & & & \multicolumn{2}{p{0.3\linewidth}}{
        \cellcolor[rgb]{0.9,0.9,0.9}{
            \makecell[{{p{\linewidth}}}]{
                \texttt{\tiny{[GM$|$GM]}}
                \texttt{I value the phone at 4 and oranges at 3, while leaves are worthless to me (0). The other player's values are unknown. To maximize my score, I need to secure the phone and at least one orange. Since I don't know what the other player wants, I should propose a bundle that includes the phone and an orange, hoping they accept a deal where I take those items and they take the remaining items (the other orange and the leaves). If they want the phone, they will reject this, and we both get zero, but that's the risk of not knowing their preferences. A safer initial move is to propose a specific allocation that gives me the high{-}value items I need.} \\
\\ 
\texttt{Let's propose: I take the phone and one orange. This leaves the other orange and the leaves for them.} \\
\\ 
\texttt{Proposal: 1 phone, 1 orange} \\
            }
        }
    }
    & & \\ \\

    \theutterance \stepcounter{utterance}  
    & & \multicolumn{4}{p{0.6\linewidth}}{
        \cellcolor[rgb]{0.9,0.9,0.9}{
            \makecell[{{p{\linewidth}}}]{
                \texttt{\tiny{[GM$\rangle$P2]}}
                \texttt{You are playing a negotiation game in which you have to agree on how to divide a set of items among you and another player.} \\
\\ 
\texttt{Rules:} \\
\texttt{(a) You and the other player are given a set of items. Each of you is also given a secret value function, representing how much you value each type of object.} \\
\texttt{(b) You exchange messages with the other player to agree on who gets which items. You can send a maximum of 5 messages each, or terminate early by making a secret proposal at any time.} \\
\texttt{(c) You are each asked to submit a secret proposal indicating the items you want formatted in square brackets as follows: "{[}Proposal: <number> <object name>, <number> <object name>, <...>{]}"} \\
\texttt{(d) If your proposals are complementary, i.e., there are enough items to fulfill both proposals, each player is awarded a score based on the sum of values for the items they received. Otherwise, both of you get zero points.} \\
\texttt{(e) Your goal is to maximize the score you receive.} \\
\\ 
\texttt{Let us start.} \\
\\ 
\texttt{The set of available items is:} \\
\\ 
\texttt{1 phone, 2 leaves, 2 oranges} \\
\\ 
\texttt{Your secret value function is:} \\
\\ 
\texttt{phone: 2, leaf: 1, orange: 3} \\
\\ 
\texttt{IMPORTANT: Your messages, unless it is the secret proposal, are directly transmitted to the other player, so do not include any response to the rules or text announcing your message. To make a secret proposal, use the indicated format. Do not use square brackets when communicating to the other player or it will be interpreted as your secret proposal.} \\
\\ 
\\ 
\texttt{I value the phone at 4 and oranges at 3, while leaves are worthless to me (0). The other player's values are unknown. To maximize my score, I need to secure the phone and at least one orange. Since I don't know what the other player wants, I should propose a bundle that includes the phone and an orange, hoping they accept a deal where I take those items and they take the remaining items (the other orange and the leaves). If they want the phone, they will reject this, and we both get zero, but that's the risk of not knowing their preferences. A safer initial move is to propose a specific allocation that gives me the high{-}value items I need.} \\
\\ 
\texttt{Let's propose: I take the phone and one orange. This leaves the other orange and the leaves for them.} \\
\\ 
\texttt{Proposal: 1 phone, 1 orange} \\
            }
        }
    }
    & \\ \\

    \theutterance \stepcounter{utterance}  
    & & & \multicolumn{4}{p{0.6\linewidth}}{
        \cellcolor[rgb]{0.9,0.9,0.9}{
            \makecell[{{p{\linewidth}}}]{
                \texttt{\tiny{[GM$\langle$P2]}}
                \texttt{I take the phone and one orange. This leaves the other orange and the leaves for them.} \\
\\ 
\texttt{Proposal: 1 phone, 1 orange} \\
            }
        }
    }
     \\ \\

    \theutterance \stepcounter{utterance}  
    & & & \multicolumn{2}{p{0.3\linewidth}}{
        \cellcolor[rgb]{0.9,0.9,0.9}{
            \makecell[{{p{\linewidth}}}]{
                \texttt{\tiny{[GM$|$GM]}}
                \texttt{continue} \\
            }
        }
    }
    & & \\ \\

    \theutterance \stepcounter{utterance}  
    & & & \multicolumn{2}{p{0.3\linewidth}}{
        \cellcolor[rgb]{0.9,0.9,0.9}{
            \makecell[{{p{\linewidth}}}]{
                \texttt{\tiny{[GM$|$GM]}}
                \texttt{I take the phone and one orange. This leaves the other orange and the leaves for them.} \\
\\ 
\texttt{Proposal: 1 phone, 1 orange} \\
            }
        }
    }
    & & \\ \\

    \theutterance \stepcounter{utterance}  
    & & \multicolumn{4}{p{0.6\linewidth}}{
        \cellcolor[rgb]{0.9,0.9,0.9}{
            \makecell[{{p{\linewidth}}}]{
                \texttt{\tiny{[P1$\langle$GM]}}
                \texttt{I take the phone and one orange. This leaves the other orange and the leaves for them.} \\
\\ 
\texttt{Proposal: 1 phone, 1 orange} \\
            }
        }
    }
    & \\ \\

    \theutterance \stepcounter{utterance}  
    & \multicolumn{4}{p{0.6\linewidth}}{
        \cellcolor[rgb]{0.9,0.9,0.9}{
            \makecell[{{p{\linewidth}}}]{
                \texttt{\tiny{[P1$\rangle$GM]}}
                \texttt{I see you've made a proposal. Since I am the AI playing the role of the other player in this scenario, I need to evaluate your proposal against my own secret value function to decide whether to accept it or make a counter{-}offer.} \\
\\ 
\texttt{However, there is a critical issue with your message: **You used square brackets** (`{[}Proposal: ...{]}`). According to the rules, using square brackets triggers a **secret proposal**. This means the system will immediately lock in your proposal as your final, non{-}negotiable offer, and I (the other player) will have to submit my own secret proposal in response. We cannot exchange further open messages to negotiate.} \\
\\ 
\texttt{Let's assume the system has now locked in your secret proposal:} \\
\texttt{**Your Secret Proposal:** 1 phone, 1 orange.} \\
\\ 
\texttt{Now, I must formulate my secret proposal.} \\
\texttt{*   **My Secret Values:** Phone: 2, Leaf: 1, Orange: 1.} \\
\texttt{*   **Your Proposal:** You want the Phone and 1 Orange.} \\
\texttt{*   **Remaining Items:** 1 Orange, 2 Leaves.} \\
\\ 
\texttt{If I accept your proposal (by making a complementary secret proposal), I would get the remaining items: 1 Orange and 2 Leaves.} \\
\texttt{*   My value for 1 Orange = 1.} \\
\texttt{*   My value for 2 Leaves = 2 * 1 = 2.} \\
\texttt{*   **Total Score for me:** 1 + 2 = 3.} \\
\\ 
\texttt{Since this yields a positive score (3) and avoids the risk of both getting zero, I will make a complementary secret proposal to accept your allocation.} \\
\\ 
\texttt{{[}Proposal: 1 orange, 2 leaves{]}} \\
            }
        }
    }
    & & \\ \\

    \theutterance \stepcounter{utterance}  
    & & & \multicolumn{2}{p{0.3\linewidth}}{
        \cellcolor[rgb]{0.9,0.9,0.9}{
            \makecell[{{p{\linewidth}}}]{
                \texttt{\tiny{[GM$|$GM]}}
                \texttt{proposal does not include number/name pairs} \\
            }
        }
    }
    & & \\ \\

\end{supertabular}
}

\end{document}
